\chapter{Verhaaltjes oplossen}\label{ch:g2:problems}

In een verhaaltje met getallen zit een vraag verstopt: optellen?
aftrekken? vermenigvuldigen? Dit hoofdstuk geeft een manier om te
kiezen, gebruikt eenvoudige strookjes om de bewerking te \emph{zien},
en leest gegevens uit tabellen en pictogrammen.

\section{De bewerking kiezen}

\begin{method}[Welke bewerking?]\label{met:g2:problems:choose}
\begin{enumerate}
  \item \emph{samenvoegen, in totaal}: optellen ($+$);
  \item \emph{wegnemen, hoeveel blijft over, hoeveel meer}: aftrekken
        ($-$);
  \item \emph{zoveel gelijke groepjes}: vermenigvuldigen ($\times$);
  \item maak snel een tekening als je twijfelt, en antwoord altijd met
        een zin.
\end{enumerate}
\end{method}

\begin{omfigure}
\begin{tikzpicture}[scale=0.6]
  % samenvoegen
  \draw[thick, fill=omDef!20] (0,0) rectangle (3,0.8);
  \draw[thick, fill=omThm!20] (3,0) rectangle (4.6,0.8);
  \node[font=\small] at (1.5,0.4) {$24$};
  \node[font=\small] at (3.8,0.4) {$15$};
  \draw[Stealth-Stealth] (0,-0.4) -- (4.6,-0.4);
  \node[below, font=\small] at (2.3,-0.5) {? in totaal: $24 + 15$};
  % vergelijken
  \begin{scope}[xshift=8cm]
  \draw[thick, fill=omDef!20] (0,0.9) rectangle (4.6,1.7);
  \node[font=\small] at (2.3,1.3) {$40$};
  \draw[thick, fill=omThm!20] (0,0) rectangle (2.8,0.8);
  \node[font=\small] at (1.4,0.4) {$28$};
  \draw[Stealth-Stealth] (2.8,0.4) -- (4.6,0.4);
  \node[below, font=\small] at (2.3,-0.2) {? meer: $40 - 28$};
  \end{scope}
\end{tikzpicture}

{\small Strookjes: twee strookjes achter elkaar betekenen een
optelling; twee strookjes naast elkaar vergeleken betekenen een
aftrekking.}
\end{omfigure}

\begin{example}[Een verhaaltje in twee stappen]\label{ex:g2:problems:twostep}
``Sam koopt $3$ pakjes van $6$ stickers en geeft er daarna $5$ weg.
Hoeveel houdt hij er over?''
\begin{enumerate}
  \item Groepjes: $3 \times 6 = 18$ stickers;
  \item weggeven: $18 - 5 = 13$.
\end{enumerate}
Sam houdt $13$ stickers over. Lees de vraag nog eens: is het verhaal
echt af? Hier wel.
\end{example}

\begin{example}[Een getal te veel]\label{ex:g2:problems:extra}
``Een boek van $210$ bladzijden\,\dots'' --- pas op, niet elk getal in
een verhaaltje heb je nodig. ``Ben heeft $8$ rode en $6$ blauwe
knikkers en is $7$ jaar oud. Hoeveel knikkers heeft hij?'' De leeftijd
is een valstrik: $8 + 6 = 14$ knikkers. Sorteer de getallen voordat je
gaat rekenen.
\end{example}

\section{Tabellen en pictogrammen}

\begin{example}[Een tabel]\label{ex:g2:problems:table}
Verkochte gebakjes op een marktkraam:
\begin{center}
\renewcommand{\arraystretch}{1.3}
\begin{tabular}{l|ccc}
 & zaterdag & zondag & totaal \\
\hline
muffins & $12$ & $9$ & $21$ \\
taartjes & $7$ & $10$ & $17$ \\
\end{tabular}
\end{center}
Aflezen: taartjes op zondag, $10$. Rekenen: muffins over de twee dagen,
$12 + 9 = 21$; op zondag $10 - 7 = 3$ taartjes meer dan op zaterdag.
\end{example}

\begin{example}[Een pictogram]\label{ex:g2:problems:pictogram}
Lees eerst de legenda: hier staat één stip voor $2$ stuks fruit.

\begin{omfigure}
\begin{tikzpicture}[scale=0.62]
  \node[left, font=\small] at (0,2) {rood};
  \foreach \i in {0,1,2,3} \fill[omThm] (0.5+\i*0.75,2) circle (0.24);
  \node[left, font=\small] at (0,1) {groen};
  \foreach \i in {0,1} \fill[omMeth] (0.5+\i*0.75,1) circle (0.24);
  \node[left, font=\small] at (0,0) {geel};
  \foreach \i in {0,1,2} \fill[omProp] (0.5+\i*0.75,0) circle (0.24);
  \node[right, font=\small] at (4.0,1) {legenda: één stip $= 2$ stuks};
\end{tikzpicture}

{\small Een pictogram: rood heeft $4$ stippen, dus $4 \times 2 = 8$
rode stuks fruit; groen $2$ stippen, $4$ stuks; geel $3$ stippen, $6$
stuks.}
\end{omfigure}
\end{example}

\section{Oefeningen}

\begin{exercise}[$\star$]\label{exo:g2:problems:1}
Noem de bewerking (reken niet uit): ``$34$ en $28$ knikkers bij elkaar
gelegd''; ``$50$ snoepjes, $12$ opgegeten''; ``$4$ dozen van $5$
appels''; ``Ana heeft $30$, Leo heeft $18$: hoeveel heeft Ana er
meer?''.
\end{exercise}

\begin{exercise}[$\star$]\label{exo:g2:problems:2}
Los op: ``Op een parkeerterrein staan $156$ auto's; er rijden $48$ weg.
Hoeveel blijven er staan?''
\end{exercise}

\begin{exercise}[$\star$]\label{exo:g2:problems:3}
Los op: ``$5$ kinderen plukken elk $4$ bloemen. Hoeveel bloemen zijn
dat in totaal?''
\end{exercise}

\begin{exercise}[$\star$]\label{exo:g2:problems:4}
Teken strookjes en los dan op: ``Tom heeft $45$ kaarten, Mia heeft er
$17$ minder. Hoeveel heeft Mia?''
\end{exercise}

\begin{exercise}[$\star$]\label{exo:g2:problems:5}
Twee stappen: ``In een krat gaan $6$ flessen. Een winkel krijgt $4$
kratten en verkoopt $9$ flessen. Hoeveel flessen blijven over?''
\end{exercise}

\begin{exercise}[$\star$]\label{exo:g2:problems:6}
Zoek het onnodige getal, streep het door en los dan op: ``Ben is $8$
jaar en heeft $23$ knikkers. Hij wint er $15$ bij. Hoeveel knikkers
heeft hij nu?''
\end{exercise}

\begin{exercise}[$\star$]\label{exo:g2:problems:7}
Gebruik de tabel van \cref{ex:g2:problems:table}: hoeveel taartjes zijn
er over de twee dagen verkocht? En hoeveel gebakjes van alle soorten op
zaterdag?
\end{exercise}

\begin{exercise}[$\star$]\label{exo:g2:problems:8}
Gebruik het pictogram van \cref{ex:g2:problems:pictogram}: hoeveel
stuks fruit zijn er in totaal? Welke kleur heeft er het meest?
\end{exercise}

\begin{exercise}[$\star$]\label{exo:g2:problems:9}
Los op: ``Een lint is $80$ cm lang. Lea knipt er $35$ cm af. Hoe lang
is het stuk dat overblijft?''
\end{exercise}

\begin{exercise}[$\star\star$]\label{exo:g2:problems:10}
``Een bakker maakt 's ochtends $5$ platen van $4$ gebakjes en 's
middags $3$ platen van $4$ gebakjes. Hoeveel gebakjes maakt hij die
dag?'' (Twee vermenigvuldigingen en dan een \omterm{def:g1:addition:def}{som} --- of tel eerst de
platen!)
\end{exercise}
